\section{情報のCIAと不審な依頼への初動}
\label{theory:research:security}
情報の機密性は許可された者に利用を限ること、完全性は情報の正確さ・完全さ、
可用性は必要なときに使えることである。
例えば設計値が書き換えられれば完全性、正しい値を必要時に取り出せなければ可用性の問題となる。
秘密を守るアクセス制限だけでは、他の二つを保証しない。\cite{saf53_ipa_cia}

方針を定め、対策を実施し、監査や監視で有効性を評価して見直す。
幹部を含む教育、事故時の初動、分析、復旧と再発防止を組織の役割へつなぐ。
見知った送信者名でも、時期や依頼内容に異常があれば、メール内の窓口に確認を求めるだけでは独立確認にならない。
既知の担当部署へ速やかに報告し、担当者が真偽と他の受信者への影響を調べる。
2017年のIPA資料は、開封率だけでなく報告から対応する組織の体制も訓練対象とする。\cite{saf53_ipa_mail}
一つの違和感は攻撃の確定証拠ではないが、被害が出るまで待つ理由にもならない。
